\documentclass[../main]{subfiles}
\begin{document}

\chapter{Segundo Principio de la Termodinámica}

\section{Introducción al Segundo Principio de la Termodinámica}

El Segundo Principio de la Termodinámica es una ley fundamental en la física que establece la dirección del flujo de calor y la imposibilidad de ciertos procesos termodinámicos. Una de las formulaciones más conocidas es la de Kelvin-Planck, que establece que es imposible construir una máquina térmica que, operando en un ciclo, no produzca otro efecto que no sea la absorción de calor de un reservorio y la realización de un trabajo.

\info{Segundo Principio de la Termodinámica}{
El Segundo Principio de la Termodinámica se puede enunciar de varias maneras, pero una de las más conocidas es la formulación de Kelvin-Planck, que establece que es imposible construir una máquina térmica que, operando en un ciclo, no produzca otro efecto que no sea la absorción de calor de un reservorio y la realización de un trabajo.
}

\subsection{Entropía}

La entropía es una medida del desorden o la aleatoriedad de un sistema. Según el Segundo Principio de la Termodinámica, la entropía de un sistema aislado tiende a aumentar con el tiempo.

\subsection{Eficiencia de una máquina térmica}

La eficiencia de una máquina térmica se define como la relación entre el trabajo realizado por la máquina y el calor absorbido de la fuente caliente. Se calcula mediante la fórmula:

\begin{equation}
    \text{Eficiencia} = \frac{\text{Trabajo}}{\text{Calor absorbido}}
\end{equation}

\subsection{Procesos reversibles e irreversibles}

Un proceso termodinámico reversible es aquel que puede revertirse completamente, dejando el sistema y su entorno en el mismo estado en el que estaban antes del proceso. Por otro lado, un proceso irreversible es aquel que no puede revertirse por sí solo.

\actividad{Responder el siguiente cuestionario}
\begin{enumerate}
    \item ¿Qué establece el Segundo Principio de la Termodinámica según la formulación de Kelvin-Planck?
    \item ¿Cuál es la relación entre el Segundo Principio de la Termodinámica y la dirección del flujo de calor?
    \item ¿Qué es una máquina térmica y cómo se relaciona con el Segundo Principio de la Termodinámica?
    \item ¿Qué es un ciclo termodinámico y cuál es su importancia en el estudio de las máquinas térmicas?
    \item ¿Qué se entiende por eficiencia de una máquina térmica y cómo se calcula?
    \item ¿Cuál es la importancia del Segundo Principio de la Termodinámica en la tecnología y la industria?
    \item ¿Qué es la entropía y cómo se relaciona con el Segundo Principio de la Termodinámica?
    \item ¿Qué son los procesos reversibles e irreversibles en el contexto del Segundo Principio de la Termodinámica?
    \item ¿Cuál es la diferencia entre una máquina térmica y un refrigerador en términos de su funcionamiento según el Segundo Principio de la Termodinámica?
    \item ¿Qué aplicaciones prácticas tiene el Segundo Principio de la Termodinámica en la vida cotidiana?
\end{enumerate}

\actividad{Resolver los siguientes problemas}
\begin{enumerate}
    \item Una máquina térmica opera entre dos reservorios a $500 \, \text{K}$ y $300 \, \text{K}$, respectivamente. Si la máquina produce un trabajo de $2000 \, \text{J}$ en cada ciclo, ¿cuánto calor absorbe del reservorio caliente en cada ciclo?
    \item Un refrigerador opera entre un ambiente a $27 \, \text{°C}$ y un espacio interno a $5 \, \text{°C}$. Si el refrigerador extrae $400 \, \text{J}$ de calor del espacio interno por cada $1000 \, \text{J}$ de trabajo que realiza, ¿cuál es la eficiencia del refrigerador?
    \item Una máquina térmica tiene una eficiencia del $40\%$ y produce $600 \, \text{J}$ de trabajo en cada ciclo. ¿Cuánto calor absorbe del reservorio caliente en cada ciclo?
    
\end{enumerate}
\end{document}